\documentclass[11pt]{article}
\usepackage[OT1]{fontenc}
\renewcommand{\rmdefault}{cmr}
\renewcommand{\sfdefault}{cmss}
\renewcommand{\ttdefault}{cmtt}
\usepackage{amsmath,amssymb,amsthm}
\usepackage[hidelinks]{hyperref}
\setlength{\textwidth}{6.1in}
\setlength{\oddsidemargin}{0.2in}
\setlength{\evensidemargin}{0.2in}
\setlength{\textheight}{9in}
\setlength{\topmargin}{-0.35in}
\setlength{\headheight}{13.6pt}
\setlength{\emergencystretch}{1.5em}
\newtheorem{theorem}{Theorem}[section]
\newtheorem{lemma}[theorem]{Lemma}
\newtheorem{proposition}[theorem]{Proposition}
\newtheorem{corollary}[theorem]{Corollary}
\theoremstyle{definition}
\newtheorem{definition}[theorem]{Definition}
\newtheorem{example}[theorem]{Example}
\newcommand{\Comp}{\operatorname{Comp}}
\newcommand{\Forb}{\operatorname{Forb}_{\rm ind}}
\newcommand{\E}{\mathsf E}
\newcommand{\C}{\mathcal C}
\newcommand{\B}{\mathcal B}
\newcommand{\OO}{\mathcal O}
\newcommand{\NN}{\operatorname{N}}
\newcommand{\dSigma}{d\text{-}\Sigma}

\title{Joint embedding for finite order forests:\\a decision theorem and a width classification}
\author{}
\date{30 October 2026\\Working manuscript}
\begin{document}
\maketitle
\begin{abstract}
Given finitely many forbidden finite order forests, can one decide whether
any two permitted forests have a permitted common extension? We answer
this question effectively by passing to comparability graphs, decomposing
a regular language of binary union/join representations into tree ideals,
and deciding inclusion between their graph images. The transfer respects existence of embeddings,
although a particular graph embedding need not preserve the order.
We also prove an explicit classification: forests of width at most $m$
have the joint embedding property exactly when $m\leq3$ (with width zero
included if the empty forest is allowed). The negative cases have finite
certificates, and the positive cases yield countable forests with the
prescribed age. The algorithm uses the effective ranked-tree ideal
decomposition of Goubault-Larrecq and Schmitz and the effective Parikh
theorem. The graph-image, inclusion, order-theoretic and width arguments
are given in full.
\end{abstract}

\section{The question and the role of the graph reduction}
An order forest is a partially ordered set in which the predecessors of
each element form a chain. A finite forest can have several minimal
elements. An embedding preserves and reflects the order. In particular,
it need not preserve a root, an immediate-predecessor relation, or a meet;
none of these is part of the language.

For a finite family $\mathcal U$ of finite order forests, put
\[
 \mathcal C(\mathcal U)=\{T:T\text{ is a finite order forest and no }
 U\in\mathcal U\text{ embeds into }T\}.
\]
A class has the \emph{joint embedding property} (JEP) if any two members
embed into a third member. Gonzalez, Rossegger and Turetsky ask which
finite forbidden antichains have this property in Question~2 of
\cite{GRT}. Here an antichain of forbidden structures is under the
embedding relation; it need not consist of poset antichains.

The question is algorithmic as well as structural. Failure of JEP can
be witnessed by two finite forests, but searching for such a pair does
not decide the positive case. The missing positive termination guarantee
comes from an effective finite ideal cover. Two issues require proof:
translating order embeddings into graph embeddings, and recognizing all
binary representations of a graph when testing the ideal cover. Related
graph JEP decision questions were studied by Carter~\cite{Carter}; the
algorithm proved here uses the independent tree-language inputs stated
in Section~3.

As a first example, forbidding a two-element antichain leaves just chains.
Two finite chains embed into a sufficiently long chain, so JEP holds.
Forbidding a five-element antichain instead allows width four. We shall
see that a balanced four-leaf binary tree and a four-leaf comb have no
common extension of width four. Thus the obstruction is branching
shape, even though the embeddings do not name branching operations.

\section{Recovering the order from its comparability graph}
Write $\Comp(T)$ for the simple graph on $T$ in which distinct vertices
are adjacent exactly when they are comparable. Write $P_4$ for the
four-vertex path and $C_4$ for the four-vertex cycle. Graph embeddings
below are induced embeddings. For a graph family $\mathcal F$,
$\Forb(\mathcal F)$ denotes the finite graphs avoiding every member of
$\mathcal F$ as an induced subgraph.

\begin{lemma}\label{lem:bridge}
The comparability graphs of finite order forests are precisely
$\Forb(P_4,C_4)$. Each such graph determines its order forest up to
isomorphism. Moreover, for finite order forests $S,T$,
\[
 S\hookrightarrow T\quad\Longleftrightarrow\quad
 \Comp(S)\hookrightarrow\Comp(T).
\]
\end{lemma}
\begin{proof}
Each connected component of a finite forest has a least element: the
root beneath a vertex is unchanged along any edge of its comparability
graph, and hence along a connected component. Its graph is obtained by
adding a universal vertex to the disjoint union of the graphs above
that root. Recursion excludes induced $P_4$ and $C_4$: each forbidden
graph is connected and neither has a universal vertex.

Conversely, a nonempty connected graph avoiding $P_4,C_4$ has a universal
vertex. Choose $v$ with maximal size of its closed neighborhood $N[v]$.
If $v$ is not universal, a shortest path to a nonneighbor has length
two, since a longer shortest path would contain an induced $P_4$.
Call the path $v,w,x$. Because $x\in N[w]\setminus N[v]$ and
$|N[v]|\geq|N[w]|$, choose $u\in N[v]\setminus N[w]$. The four distinct
vertices $u,v,w,x$ induce $P_4$ when $ux$ is absent and $C_4$ when it is
present. Both are impossible. Remove a universal vertex, recursively
represent the remaining components, and place the removed vertex below
them. Represent a disconnected graph by a disjoint union.

For uniqueness, all universal vertices in a connected representing
forest form an initial chain. They are pairwise comparable. If $z\leq u$
and $u$ is universal, then $z$ is comparable with every vertex: use
transitivity above $u$, and linearity of the predecessors of $u$ below
it. Thus $z$ is universal too. Every nonuniversal vertex lies above the
whole universal chain. Deleting that chain leaves the uniquely specified
graph components. Induction reconstructs their orders up to isomorphism.
Permuting the names of the universal vertices does not change that
isomorphism type.

An order embedding gives an induced graph embedding. For the converse,
give the image of a graph embedding the order induced from $T$. It is
a forest whose comparability graph is isomorphic to $\Comp(S)$.
Uniqueness gives an order isomorphism from $S$ to that image.
\end{proof}

The last argument asserts existence of an order embedding, not that
the original graph map is an order embedding. For a two-element chain,
the transposition is a graph automorphism but reverses the order.
That distinction is harmless for JEP and essential for interpreting
the reduction correctly. These graphs are commonly called trivially
perfect graphs; the proof above supplies all properties needed here.

\section{Deciding the forbidden-family question}

The comparability-graph correspondence reduces the question to a hereditary
class of graphs. A graph can have many binary union/join representations,
and those representations need not embed into a common tree even when their
graphs have a common host. We therefore decompose a language of representations
into directed pieces and compare their \emph{graph images}. This section
gives the effective steps, including the passage between different binary
representations of the same graph.

We use two standard effective results. First, for a regular language of
ranked trees, one can compute its homeomorphic downward closure as a finite
union of regular ideals. The precise input is the construction of
Goubault-Larrecq and Schmitz~\cite{GLS}, Section~5: Proposition~18 constructs
simple tree regular expressions; Fact~19 and Lemma~20 give a terminating
normalization; Theorem~23 identifies the resulting tree products as ideals.
Second, the Parikh image of a context-free language has an effective
Presburger description. One explicit construction is given by Esparza,
Ganty, Kiefer and Luttenberger~\cite{Parikh}, Theorem~1.1 and Section~3.
Presburger arithmetic is first-order arithmetic with addition and order on
the natural numbers; its sentences are decidable. These inputs concern tree
languages and finite grammars. The graph-containment and JEP transfers below
are proved here.

\subsection{Binary representations and ideals}
Use the ranked alphabet $\{a:0,U:2,J:2\}$. Each occurrence of $a$ represents
one vertex. The operations $U$ and $J$ represent disjoint union and join,
respectively. Write $\E(t)$ for the graph represented by a term $t$.
Children are ordered in a term, although graph evaluation forgets their
order. The empty graph is treated separately.

The ordered homeomorphic relation $\preceq$ on these terms has the following
recursive definition. A source can embed into either child of a target;
alternatively, equal root symbols can be coupled with childwise embeddings.
Thus $a\preceq a$, and for $o\in\{U,J\}$ the coupling clause is
\[
 o(s_1,s_2)\preceq o(t_1,t_2)
 \quad\hbox{if }s_1\preceq t_1\hbox{ and }s_2\preceq t_2.
\]
This is exactly the ranked relation in Example~7 of~\cite{GLS}.
It preserves the root operation at each coupled node. An \emph{ideal}
is a nonempty downward-closed directed set under the specified relation.

\begin{lemma}\label{lem:parse-pruning}
If $s\preceq t$, then $\E(s)$ embeds as an induced subgraph of $\E(t)$.
For every nonempty subset $B$ of the leaves of $t$, there is a term
$t|B\preceq t$ representing exactly $\E(t)[B]$.
\end{lemma}
\begin{proof}
For the first assertion, diving includes a child graph as an induced
subgraph of its parent. Coupling combines two child embeddings on disjoint
vertex sets: the common operation puts no edges, or all edges, between
those sets on both sides. Induction proves the assertion.

For the second, delete leaves outside $B$ recursively. At a node discard
empty child results; return the surviving child if there is only one,
and retain the original binary operation if both survive. Suppressing a
unary node keeps the surviving child's own operation. Restriction commutes
with both graph operations, which proves the graph identity. The dive and
coupling clauses prove $t|B\preceq t$ by induction. Adjacent equal
operations need not be flattened: the result is a binary representation.
\end{proof}

Let $\mathcal F$ be a finite forbidden graph family containing $P_4,C_4$,
and put $\mathcal K=\Forb(\mathcal F)$. Every nonempty member of
$\mathcal K$ has a binary representation: the construction in
Lemma~\ref{lem:bridge} recursively takes disjoint unions and adds a
universal vertex by joining with $a$.

\begin{lemma}\label{lem:parse-regular}
The language $L=\{t:\E(t)\in\mathcal K\}$ is effectively regular and
downward closed under $\preceq$.
\end{lemma}
\begin{proof}
Choose $n\geq1$ at least as large as every forbidden graph. A state records
the set of induced graph isomorphism types with at most $n$ vertices present
in the represented graph, including the empty type. There are finitely many
such types, and canonical representatives and isomorphism tests are finite
computations. The leaf state is $\{\varnothing,K_1\}$. For either operation
$o$, the transition on states $S,T$ is
\[
 \{o(G,H):G\in S,\ H\in T,\ |G|+|H|\leq n\},
\]
with graph types identified up to isomorphism. The empty graph is an
identity for both operations. Every induced restriction of an operation
splits into restrictions of its two arguments, so these transitions compute
the stated profile exactly. Accept states containing no member of
$\mathcal F$. This constructs the required finite automaton. Heredity
and Lemma~\ref{lem:parse-pruning} give downward closure.
\end{proof}

Apply the effective ideal-decomposition theorem to obtain
\[
                     L=I_1\cup\cdots\cup I_r,                 \tag{1}
\]
where the $I_i$ are regular term ideals with effectively available automata.
If a forbidden graph is empty, $\mathcal K$ is empty and has vacuous JEP.
Otherwise it contains the empty graph. If $L$ is empty, $\mathcal K$
contains only that graph and has JEP too. These cases are decidable;
henceforth $r\geq1$.

\begin{lemma}\label{lem:image-cover}
Put $D_i=\{\varnothing\}\cup\{\E(t):t\in I_i\}$, on isomorphism types.
Each $D_i$ is a hereditary graph class with JEP, their union is
$\mathcal K$, and
\[
 \mathcal K\text{ has JEP}\quad\Longleftrightarrow\quad
                     \mathcal K=D_i\text{ for some }i.        \tag{2}
\]
\end{lemma}
\begin{proof}
Containment and covering follow from (1) and the representation of every
nonempty graph in $\mathcal K$. For heredity take a representation in
$I_i$ and apply the pruning part of Lemma~\ref{lem:parse-pruning}; the
result stays in $I_i$. For JEP take two representing terms in $I_i$,
use its directedness, and evaluate a common term host. Empty graphs cause
no difficulty.

Equality with one $D_i$ implies JEP. Conversely, suppose $\mathcal K$
has JEP but each $D_i$ is proper. Choose $G_i\in\mathcal K\setminus D_i$.
Iterating JEP supplies one $H\in\mathcal K$ containing all the finitely
many $G_i$. The cover puts $H$ in some $D_j$; heredity then puts $G_j$
there, a contradiction.
\end{proof}

The remaining task is to decide the finitely many inclusions
$\mathcal K\subseteq D_i$. Individual binary representations cannot be
used for this test, since a graph may have another representation accepted
by $I_i$. We now give the required exact test.

\subsection{Recognizing graph images of regular term languages}
A \emph{reduced cotree} has leaves $a$ and unordered internal nodes
$U,J$ of arity at least two, with adjacent internal operations unequal.
Flattening adjacent equal operations turns a binary representation into
such a cotree without changing its graph.

\begin{lemma}\label{lem:canonical-cotree}
Every graph with a union/join representation has a unique reduced cotree
up to isomorphism. All its binary representations are obtained by
independently representing each cotree child, then choosing a binary
association and a permutation of the children at each operation node.
\end{lemma}
\begin{proof}
At a reduced $U$ root each child graph is connected: a child is a leaf or
has a join root. Its children are therefore precisely the graph's connected
components. At a $J$ root the graph is connected and its children are
precisely the connected components of the complement. These decompositions
recover the root operation and its child vertex sets from the graph.
Induction recovers every child cotree; zero and one vertices are immediate.
This proves uniqueness. Flattening any binary representation yields this
cotree, and reversing each flattening uses precisely binary association
and permutation within a single operation. This proves the last assertion.
\end{proof}

Let $M$ be a deterministic complete binary tree automaton with finite
state set $Q$, leaf state $q_a$, transitions $\delta_U,\delta_J$, and
accepting set $A_M$. For a reduced cotree $T$ define
\[
 S_M(T)=\{q\in Q:\text{some binary representation of }T
                         \text{ has state }q\}.              \tag{3}
\]
Storing the full subset accounts for all possible representations.

\begin{lemma}\label{lem:parikh-transition}
For each $o\in\{U,J\}$ and $q\in Q$ one can compute a Presburger
predicate $P_{o,q}$ on the counts of child state subsets such that, if
an $o$ node has children $T_1,\ldots,T_k$, then
\[
 q\in S_M(T)\quad\Longleftrightarrow\quad
 P_{o,q}\bigl((c_B)_{B\subseteq Q}\bigr),\qquad
 c_B=|\{j:S_M(T_j)=B\}|.                                   \tag{4}
\]
Equal child occurrences are counted with their full multiplicity.
\end{lemma}
\begin{proof}
Use nonterminals $N_q$ for $q\in Q$ and one terminal $b_B$ for every
subset $B\subseteq Q$. For the operation $o$ form the finite grammar
\[
 N_{\delta_o(r,s)}\longrightarrow N_rN_s\quad(r,s\in Q),
 \qquad N_q\longrightarrow b_B\quad(q\in B).
\]
Compute the Parikh predicate for start nonterminal $N_q$. A binary
representation of the node consists of a binary $o$ skeleton, with each
child represented at a chosen state belonging to its subset. Reading
that skeleton gives a derivation with the stated terminal counts.

Conversely a derivation with those counts supplies an $o$ skeleton and
one selected state for each terminal occurrence. Match occurrences to
the actual children of the same subset type. For each occurrence choose
a child representation attaining its selected state, and assemble the
representations in the skeleton's order. This is a representation of
$T$ with root state $q$. Lemma~\ref{lem:canonical-cotree} ensures that
these are exactly the allowed representations. The effective Parikh
theorem therefore proves both the correspondence and effectivity.
\end{proof}

The resulting deterministic cotree automaton has states $(\ell,S)$,
where $\ell\in\{a,U,J\}$ records the root kind and $S\subseteq Q$.
Its leaf state is $(a,\{q_a\})$. At an $o$ node, transition to a
specified subset $S$ is expressed by
\[
 \bigwedge_{q\in S}P_{o,q}(c)\ \land\
 \bigwedge_{q\notin S}\neg P_{o,q}(c).                        \tag{5}
\]
Require at least two children and exclude children with operation $o$.
The counts $c_B$ are sums of the counts for the permitted child root
kinds with subset $B$. These are linear expressions, so the horizontal
transitions remain Presburger formulas. Accept when $S\cap A_M$ is
nonempty. Induction using (4) shows that this automaton accepts exactly
the graphs represented by terms accepted by $M$, in their canonical
cotree form.

\begin{lemma}\label{lem:presburger-inclusion}
Inclusion of languages recognized by these deterministic cotree
automata is decidable.
\end{lemma}
\begin{proof}
Work throughout on the domain of reduced cotrees. Complement changes
the accepting set. For an intersection, use pairs of states; replace
each component child-state count by the sum of pair-state counts
projecting to it, and conjoin the transition formulas. Root kinds,
arity and alternation agree in the product. These are finite
Presburger constructions.

For emptiness, start with the actual leaf states and compute a least
fixed point of productive states. Given the current finite set $R$,
add an internal state if its transition formula has a nonnegative
count vector of total at least two, with count zero on every child
type outside $R$. Presburger decision settles each such question.
There are finitely many states, so the process terminates.

Each newly added state has a finite cotree representative: take the
required number of copies of representatives for its already productive
child states. Conversely induction on cotree height puts every state
of an actual finite cotree into the fixed point. Thus a language is
nonempty exactly when an accepting state is productive. Apply this to
the first language intersected with the complement of the second.
\end{proof}

\subsection{The resulting algorithm}
For the automata of $L$ and each $I_i$, construct the cotree automata
(3)--(5). They recognize $\mathcal K$ and $D_i$, respectively, with
the common empty graph omitted. Lemma~\ref{lem:presburger-inclusion}
decides $\mathcal K\subseteq D_i$. Since $D_i\subseteq\mathcal K$
already, this tests equality. By (2), return yes if some equality
holds and no otherwise. Every stage terminates, by the cited effective
theorems or the explicit finite-state procedure just proved.

For example, the independent graphs with at most four vertices have
JEP. Their binary representations include a four-leaf balanced tree and
a four-leaf comb, which cannot have a common binary host with at most
four leaves under $\preceq$. Our algorithm decomposes the representation
language instead of testing its JEP. In this finite example the maximal
binary representations generate principal term ideals. Each such ideal
has graph image exactly the independent graphs of sizes one through
four, by pruning and leaf counts. Their graph images therefore give
the correct affirmative verdict.

\begin{theorem}\label{thm:decision}
There is an algorithm which, given any finite family $\mathcal U$ of
finite order forests, decides whether $\mathcal C(\mathcal U)$ has JEP.
More precisely, set
\[
 \mathcal F=\{P_4,C_4\}\cup\{\Comp(U):U\in\mathcal U\}.
\]
Then $\mathcal C(\mathcal U)$ has JEP if and only if $\Forb(\mathcal F)$
has JEP.
\end{theorem}
\begin{proof}
Lemma~\ref{lem:bridge} gives an effective correspondence on isomorphism
types between finite forests and $\Forb(P_4,C_4)$, preserving and
reflecting existence of embeddings. Thus avoiding $\mathcal U$
corresponds exactly to avoiding the displayed graph family, and common
hosts correspond in both directions. Compute this finite graph list
and run the terminating graph procedure established above.
\end{proof}

An input family need not be an antichain: remove redundant forbidden
members by finite embedding tests. The algorithm supplies an effective
criterion, with no claim that its large intermediate automata are
practical or that the decision procedure has been implemented.


\begin{proposition}[Finite certificates for negative instances]
\label{prop:certificate}
If $\mathcal C(\mathcal U)$ fails JEP, one can compute a pair $F,G$ in
the class and a finite exhaustive certificate that they have no common
extension in the class.
\end{proposition}
\begin{proof}
Any common host can be replaced by its induced suborder on the union
of the two embedding images. Heredity preserves membership in the
class, and the resulting host has at most $|F|+|G|$ vertices. For a
fixed pair it therefore suffices to enumerate all forests of size at
most that bound and test their membership and all candidate embeddings.
These are finite calculations on finite order relations.

After a negative verdict from Theorem~\ref{thm:decision}, enumerate
permitted pairs and perform this finite test for each. A bad pair must
eventually appear. The pair together with the exhaustive host test is
a certificate verifiable by finite host enumeration alone.
\end{proof}

The bound is a bound on the size of a host for a fixed pair. It is not
a uniform bound on the size of the first bad pair as a function of the
forbidden family. Without the total decision procedure the same search would
semidecide only failure of JEP.

\section{A sharp classification by width}
The width of a finite poset is the largest size of an antichain. For a
finite forest it is the number of maximal elements, which we call leaves.
Indeed, each element of an antichain extends to a leaf, and incomparable
elements have disjoint upper cones: a common upper bound would have
incomparable predecessors.

\begin{theorem}\label{thm:width}
For $m\geq1$, the class of finite order forests of width at most $m$
has JEP if and only if $m\leq3$.
\end{theorem}
\begin{proof}
First suppose $m\leq3$. A forest of width at most $m$ embeds as an
induced suborder of a rooted binary tree with exactly $m$ leaves.
Add a root below all components; replace any multifurcation by binary
branching through inserted vertices; and, if necessary, attach extra
leaves. These operations preserve the old order and need at most $m$
leaves. Suppressing unary vertices now gives a full binary tree.
For each of $m=1,2,3$ there is exactly one unordered full binary tree
with $m$ leaves. A sufficiently long subdivision of each edge of this
tree, together with a sufficiently long initial stem, therefore contains
any prescribed pair of the original forests. Choose on each segment
enough vertices for both individual embeddings; they need not have
disjoint images. This common host still has width $m$.

Let $m\geq4$. Compare a binary comb with $m$ leaves, whose root split
is $1+(m-1)$, with a full binary tree whose root split is $2+(m-2)$.
For a vertex $v$ let $L(v)$ be its set of descendant leaves. Distinct
leaf clusters form a laminar family: two are nested or disjoint.

Suppose a full binary $m$-leaf tree embeds into a host $H$ of width
at most $m$. Its leaf images are incomparable; choose a host leaf above
each one. There must be exactly $m$ host leaves, and each source-leaf
image has only its chosen host leaf above it. For a source vertex $v$,
the host leaves above its image are precisely the matched leaves in
$L(v)$. To check the nontrivial inclusion, any host leaf is above the
image of a unique source leaf $\ell$. If it is also above the image of
$v$, foresthood makes those two images comparable. Order reflection
then makes $v$ and $\ell$ comparable, so $v\leq\ell$.

A full binary tree has $2m-1$ distinct nonempty clusters. This family
is maximal laminar on the $m$ leaves. Indeed, a set compatible with
both children of every binary split must either be contained in one
child or contain that entire child; descending recursively shows that
it is already a cluster. The host's cluster family contains the
source's and is laminar, so the two families coincide. Thus a common
host for the two chosen sources would identify their cluster families
under a leaf bijection. In the comb all nonsingleton proper clusters
are nested. The other source has two disjoint nonsingleton root-child
clusters. This is impossible.
\end{proof}

If $A_k$ is the $k$-element poset antichain, then
\[
 \mathcal C(\{A_k\})\text{ has JEP}\quad\Longleftrightarrow\quad k\leq4
 \qquad(k\geq2).
\]
For $m=4$ the two negative witnesses each have seven vertices: the
balanced binary tree and the four-leaf comb. Their embeddings were not
assumed to preserve roots or meets. Cluster preservation was derived
from the saturated width bound. If empty forests are allowed, the
width-zero class consists only of the empty forest and has JEP too.

\section{Countable realizations and scope}
The age of a countable forest is the class of isomorphism types of its
finite induced suborders. The following standard age construction
explains one use of JEP without imposing amalgamation or homogeneity.

\begin{proposition}\label{prop:age}
A nonempty hereditary class $\mathcal K$ of finite forests is the age
of a finite or countable forest if and only if it has JEP. If it contains
forests of unbounded size, the realizing forest is countably infinite.
\end{proposition}
\begin{proof}
Any two finite suborders of one forest are contained in their finite
union, proving necessity. Conversely enumerate the members of
$\mathcal K$ up to isomorphism. Use JEP repeatedly to obtain a chain
of members, with each enumerated member embedding into a later stage.
Rename each successive host so that its chosen copy of the preceding
stage is an actual induced suborder. The union is a forest: every
finite violation of the order or predecessor-chain axioms would already
occur at a finite stage. Every finite suborder of the union is contained
in a stage, and hence belongs to $\mathcal K$ by heredity. Every member
of $\mathcal K$ embeds into some stage. The union therefore has exactly
the required age. Unbounded size forces the union to be infinite.
\end{proof}

Together with Theorem~\ref{thm:decision}, this decides which finite
forbidden-family forest classes are ages. The empty class is treated
separately: under the usual vacuous convention it has JEP, but is the
age of no structure. Whether the empty forest is admitted as a structure
is another convention that can be settled by a finite input check.

The source motivation in \cite{GRT} concerns Scott analysis below the
Vaught ordinal. The present paper answers its finite-forest Question~2
by an effective criterion, and gives a direct structural answer for
antichain obstructions. It does not establish a model-counting theorem
for arbitrary complete first-order theories. Nor does the age
construction alone assign a Scott rank to its union. Such conclusions
require the additional hypotheses of the relevant results in
\cite{GRT}.

\paragraph{Preparation and status.}
This manuscript was assembled with AI assistance from a research
programme directed by the repository owner. The effective ideal-decomposition
and Parikh theorems are credited as published inputs. The graph-image and
inclusion lemmas, forest transfer and width proofs above are offered for
review; no implementation or formal verification of the full decision
algorithm is claimed. No historical priority
claim is made for those deductions. The repository contains the full
manuscript and its build instructions; it requires no private project
notes or conventions.

\begin{thebibliography}{9}
\bibitem{Carter} D.~Carter, \emph{On the joint embedding property for
cographs and trees}, arXiv:2409.06127v1 (2024).
\url{https://arxiv.org/abs/2409.06127v1}.
\bibitem{GLS} J.~Goubault-Larrecq and S.~Schmitz,
\emph{Deciding Piecewise Testable Separability for Regular Tree Languages},
ICALP 2016, LIPIcs~55, 97:1--97:15, Example~7 and Section~5.
\url{https://doi.org/10.4230/LIPIcs.ICALP.2016.97}.
\bibitem{Parikh} J.~Esparza, P.~Ganty, S.~Kiefer and M.~Luttenberger,
\emph{Parikh's Theorem: A simple and direct automaton construction},
arXiv:1006.3825v3 (2011), Theorem~1.1 and Section~3.
\url{https://arxiv.org/abs/1006.3825v3}.
\bibitem{GRT} D.~Gonzalez, D.~Rossegger and D.~Turetsky,
\emph{Scott Analysis below the Vaught Ordinal}, arXiv:2606.15205v2
(2026), Question~2.
\url{https://arxiv.org/abs/2606.15205v2}.
\end{thebibliography}
\end{document}
